\section{拒稿诊断与本报告的任务}

\subsection{拒稿意见的结构}

2026 年 9 月，\sieve{} 期刊版（投稿 Information Processing \& Management，2026-09-24 版）收到编辑直接拒稿（desk rejection）。拒稿信的核心意见只有两句：其一，论文缺乏足够的新颖性（lack of sufficient novelty）；其二，所报告的研究结果过于不成熟，不足以支撑发表（the research results reported are too premature for publication）。与此同时，编辑明确确认论文主题落在期刊的范围之内（falls within the aim and scope of this journal）。这三句话合在一起构成了一份罕见的、高度集中的诊断：期刊不质疑问题的重要性，也不质疑写作与格式，而是把全部否决理由压在两个具体的科学维度上——贡献的增量的真实性，以及证据链的成熟度。

这一诊断与本论文在 2025--2026 年文献环境中所处的位置直接相关。以 2023 年 LLMLingua 与 Selective Context 为起点，训练无关（training-free）的提示与上下文压缩已经从一个空白角落变成了一个快速拥挤的赛道：LongLLMLingua 加入了问题条件化与位置相关的机制，LLMLingua-2 引入了数据蒸馏的token分类器，Perception Compressor（NAACL 2025 Findings）明确以 training-free 为主张并纳入问题引导与位置分配，2026 年出版的 QASPC 则直接以 question-aware、sentence-level、training-free 三个属性作为标题。在这样的环境下，一篇论文如果把自己定位为``又一个 training-free 压缩器''，其新颖性主张必然与已有工作大面积重叠。拒稿意见正是这种重叠的反映。

\subsection{本报告的任务与方法论约束}

本报告是两阶段重构计划的第一阶段，对应拒稿意见的第一句话。任务只有一个：以可审计的方式逐篇查证 2023--2026 年的相关文献，回答``\sieve{} 相对于现有工作到底新增了什么''。如果这个问题回答不出来，正确的后续动作是修改方法本身而不是继续润色表述；如果能够回答，则答案将成为论文 v2 版本全部叙事（标题、摘要、引言、贡献列表、相关工作）的重建基础。

审计遵循三个方法论约束，均来自论文写作技能库（\texttt{20-ml-paper-writing}）中的引用规范。第一，绝不凭记忆生成文献事实：每一篇被审计文献的机制描述都必须能追溯到检索到的原文证据（摘要、全文或官方代码仓库），不允许``我记得这篇论文是\dots\dots''式的判断进入结论。第二，无法验证的单元格如实标注：当某篇文献的全文不可得、某个机制细节无法从公开渠道核实时，相应判定标注为``未验证''并说明原因，而不是猜测填充。第三，区分证据等级：全文与官方代码直接确认的判定记为 A 级，多源摘要片段交叉确认的记为 B 级，仅凭单一来源或无法核实的记为 C 级。第 \ref{sec:cards} 节的每张方法卡片末尾都标注了证据等级与来源清单，附录 \ref{sec:appendix} 给出了全部原始链接与检索日期。

第二阶段的任务——回应``结果过于不成熟''的证据补强——由配套的《实验补强计划》单独承担，两份文件共同构成对拒稿意见的完整回应。需要强调的是，两份文件之间存在依赖关系：实验补强的设计必须服务于新颖性重构后所声明的贡献，而不是盲目扩大实验规模。先定``新在哪里''，再定``凭什么信''，这个顺序不可颠倒。
